\section{Integral topology of the marked split quintic}\label{sec:integral-topology}
Choose independent linear forms $l_1,l_2$ and quartics $q_1,q_2$ on
$\PP^4$. Put $\Pi=\{l_1=l_2=0\}$ and assume that the restricted
quartics meet transversely in sixteen points, that
$X_0=\{l_1q_2-l_2q_1=0\}$ is smooth away from $\Pi$, and that
\[
 Y=\{ul_1+vl_2=0,\ uq_1+vq_2=0\}
       \subset\PP^4\times\PP^1
\]
is smooth. These are nonempty open conditions. Choose a quintic $G$
nonzero at the sixteen nodes and let
$X_t=\{l_1q_2-l_2q_1+tG=0\}$, for sufficiently small $t$.
The projective construction, the hyperplane markings $H,L$, and the
chosen smoothing ray are the global inputs. The complement ring and
attachment maps below are calculated, rather than assumed.

Let $U$ be the common compact exterior, $q:Y\to X_0$ the projection,
and $s_t:X_t\to X_0$ the continuous Milnor collapse along the ray.
Mark the degree-two link generators by the exceptional rulings.
Write
\[
 \Lambda=\ZZ^{16},\qquad \mathbf1=(1,\ldots,1),\qquad
 \Lambda_0=\ker(\Sigma:\Lambda\to\ZZ),\qquad W=\Lambda/\ZZ\mathbf1.
\]
Here $\Sigma$ is coordinate summation. Let $J_{15}$ be the
$15\times15$ all-ones matrix.

\begin{theorem}[The marked split-quintic transition]\label{thm:integrated}
The preceding construction has the following compatible properties.
\begin{enumerate}
\item The smooth threefolds $Y$ and $X_t$ are simply connected and
have torsion-free integral cohomology, with the full rings and Chern
classes in Table~\ref{tab:integral-rings}. Their third Betti numbers
are $174$ and $204$; their degree-three products are unimodular
alternating pairings times the point class. In both cases
$H^0=H^6=\ZZ$ and $H^1=H^5=0$. The integral groups of $U$ and $X_0$ are free,
with Betti vectors
\[
 (1,0,2,189,1,15,0),\qquad (1,0,1,189,2,0,1),
\]
respectively. The central ring and both specialization pullbacks are
given by \eqref{eq:central-ring} and the degree-three description in
Section~\ref{sec:specialization}; in particular
$\im s_t^*=\ker(T^*-I)$ in integral degree-three cohomology,
where $T$ is monodromy about a positive simple loop.

\item The degree-two boundary restriction is
$r(hH+\ell L)=\ell\mathbf1$. It occurs in the integral exact sequence
\begin{equation}\label{eq:split-exact-sequence}
 0\longrightarrow\ZZ H\longrightarrow\ZZ H\oplus\ZZ L
 \xrightarrow{\ r\ }\Lambda\xrightarrow{\ \partial\ }
 H^3(X_t;\ZZ)\longrightarrow H^3(U;\ZZ)\longrightarrow0.
\end{equation}
The classes $\partial e_i$ are the left Poincar\'e duals of the
vanishing spheres. They generate the primitive isotropic lattice $W$,
with the single relation $\sum_i\partial e_i=0$. Under Poincar\'e
duality the nonzero monodromy block is $\pm(I_{15}+J_{15})$, whose
Smith form is $\operatorname{diag}(1^{14},16)$. Thus the monodromy
image has index sixteen in the primitive vanishing lattice.

\item Use the actual Ricci-flat class $H+s^2L$ on $Y$ and class $H$
on $X_t$, where $s,t>0$ are independent branch parameters. Both
branches converge to the same metric completion $(X_0,d_{g_0})$,
homeomorphic to $X_0$, with the local rescaled resolved-conifold and
Stenzel limits. Every exceptional area on $Y$ equals $s^2$.
The same class $L$ has period one at every link, and its central
primitive harmonic representative $\alpha$ satisfies
\begin{equation}\label{eq:split-compatibility}
 [\alpha]=L-\tfrac45H,\qquad
 \|\alpha\|_{L^2(g_0)}^2=\tfrac{16}{5},\qquad
 \lim_{t\downarrow0}|\log t|\lambda_2^+(t)
 =\frac{2\pi\|r(L)\|^2}{16/5}=10\pi.
\end{equation}
The eigenvalue has the simplicity, type, gap and one-sided upper-error
properties of Theorem~\ref{thm:main}. The resolution and smoothing
limits require no relation between $s$ and $t$.

\item The tangent bundles on the smooth spaces are stable and their
Ricci-flat tangent connections are Hermitian--Einstein. The relative
tangent sheaf of the smoothing is flat, with central fibre $E_0$ and
\[
 0\longrightarrow E_0\longrightarrow T_{X_0}
 \longrightarrow\bigoplus_{i=1}^{16}\CC_{p_i}\longrightarrow0,
 \qquad E_0^{**}=q_*T_Y=T_{X_0},\qquad R^jq_*T_Y=0\ (j>0).
\]
These extensions agree on the regular locus. Their actual tangent
connections have the local limits and curvature-energy bounds in
Proposition~\ref{prop:tangent-metrics}.
\end{enumerate}
\end{theorem}


\subsection{The full integral rings}
Let $H,L$ be the hyperplane classes of $\PP^4,\PP^1$ restricted to $Y$.
Lefschetz gives $H_2(Y;\ZZ)=\ZZ^2$ and $H_2(X_t;\ZZ)=\ZZ$.
Together with $H_1=0$, universal coefficients and integral Poincar\'e
duality imply torsion-freeness in every degree. The cubic intersections are \eqref{eq:split-cubic}.
Take $A,B\in H^4(Y;\ZZ)$ dual to $H,L$, and let
$a_4\in H^4(X_t;\ZZ)$ satisfy $\int_{X_t}Ha_4=1$. Then:

\begin{table}[tb]
\centering
\caption{Integral cohomology and characteristic data of the compact fibres.}
\label{tab:integral-rings}
\begin{adjustbox}{max width=\textwidth,center}
\begin{tabular}{ll@{\hspace{1.5em}}l}
\toprule
 & $Y$ & $X_t$\\
\midrule
$H^2$ & $\ZZ H\oplus\ZZ L$ & $\ZZ H$\\
$H^4$ & $\ZZ A\oplus\ZZ B$ & $\ZZ a_4$\\
Degree-two products & $H^2=5A+4B$, $HL=4A$, $L^2=0$ & $H^2=5a_4$\\
Top products & $HA=LB=[\mathrm{pt}]$, $HB=LA=0$ & $Ha_4=[\mathrm{pt}]$\\
$\operatorname{rank}H^3$ & $174$ & $204$\\
$c_2$ & $50A+24B$ & $50a_4$\\
Euler characteristic & $-168$ & $-200$\\
\bottomrule
\end{tabular}
\end{adjustbox}
\end{table}

To verify the characteristic numbers directly, use
\[
 c(TY)=\frac{(1+H)^5(1+L)^2}{(1+H+L)(1+4H+L)}.
\]
This gives $c_2(Y)=6H^2+5HL$ and
$c_3(Y)=-20H^3-17H^2L$. For a quintic,
$c_2=10H^2$ and $c_3=-40H^3$. Since $H^1=H^5=0$, the Euler
numbers give the displayed third Betti numbers. The degree-three product
on each compact fibre is its unimodular alternating pairing times the
point class. With $H^0=H^6=\ZZ$, these statements specify the full rings
up to symplectic basis choices. The product subgroup generated by
$H^2$ and $HL$ in $H^4(Y)$ has index
$|\det\left(\begin{smallmatrix}5&4\\4&0\end{smallmatrix}\right)|=16$;
it is not an integral basis.

Every exceptional curve has degrees $(H\cdot C_i,L\cdot C_i)=(0,1)$.
They therefore represent the same primitive homology class $c_Y$, with
Poincar\'e dual $B$. Primitivity follows from the period of $L$.

\subsection{The common exterior and all cohomological attachments}
Remove disjoint tubular neighborhoods of the exceptional curves, obtaining
a compact exterior $U$. Its boundary components $\mathcal L_i$ are
$S^2\times S^3$. Define
\[
 \Lambda=\ZZ^{16},\quad \mathbf1=(1,\ldots,1),\quad
 \Lambda_0=\ker(\Sigma:\Lambda\to\ZZ),\quad W=\Lambda/\ZZ\mathbf1.
\]
The homological Thom sequence uses
$H_j(Y,U)=H_{j-4}(\coprod C_i)$.
The maps $H_4(Y)\to\Lambda$ and $H_6(Y)\to\Lambda$ have primitive
diagonal image: the former is exceptional intersection, witnessed by
$L\cdot C_i=1$, and the latter is the oriented diagonal fundamental
class. Thus $U$ has free integral homology and cohomology, with Betti
vector
\begin{equation}\label{eq:exterior-betti}
 (1,0,2,189,1,15,0),
\end{equation}
and exact sequences
\begin{equation}\label{eq:exterior-homology}
 0\to W\to H_3(U)\to H_3(Y)\to0,
 \qquad 0\to H_3(U)\to H_3(X_t)\to \Lambda_0\to0.
\end{equation}
The second comes from the Milnor fillings, whose cores are three-spheres.
The vanishing spheres span $W$, with relation
$\delta_1+\cdots+\delta_{16}=0$. Both successive quotients in
\eqref{eq:exterior-homology} are free; hence $W$ embeds primitively in
$H_3(X_t)$, and is isotropic because the spheres are disjoint.

In cohomology the Thom sequence kills exactly $\ZZ B$ in degree four.
Writing $\bar A=A|_U$, the even products are
\[
 H^2=5\bar A,\qquad HL=4\bar A,\qquad L^2=0.
\]
Orient each link as a boundary component of $U$ and choose generators
$\alpha_i\in H^2(\mathcal L_i)$ and $\beta_i\in H^3(\mathcal L_i)$
with $L|_{\mathcal L_i}=\alpha_i$ and $\int\alpha_i\beta_i=1$.
The symbol $\beta_i$ here has degree three; the harmonic degree-two
form remains $\vartheta_i$. The degree-three restriction is
\begin{equation}\label{eq:exterior-cohomology}
 0\to H^3(Y)\to H^3(U)\xrightarrow{r_U}\Lambda_0\to0,
\end{equation}
and degree-five restriction is the inclusion $H^5(U)=\Lambda_0\hookrightarrow\Lambda$.
Boundary multiplicativity and that injectivity imply, for $y\in H^3(U)$,
\[
 Hy=0,\qquad Ly=r_U(y)\in \Lambda_0=H^5(U).
\]
All degree-six products vanish. These formulas give the entire exterior
ring, without supplying an unspecified multiplication as additional input.

For completeness, use the Mayer--Vietoris difference $j_U^*-j_F^*$.
Let $x_i$ be the degree-two generator of the resolution filling and $z_i$
the degree-three generator of the smoothing filling, restricting to
$\alpha_i,\beta_i$. The nonzero difference maps are:
\begin{center}
\begin{adjustbox}{max width=\textwidth,center}
\begin{tabular}{cll}
\toprule
Degree & Resolution & Smoothing\\
\midrule
$0$ & $c_0\mathbf1-(c_i)$ & $c_0\mathbf1-(c_i)$\\
$2$ & $\ell\mathbf1-(x_i)$ & $\ell\mathbf1$\\
$3$ & $r_U(y)$ & $r_U(y)-(z_i)$\\
$5$ & $\Lambda_0\hookrightarrow\Lambda$ & $\Lambda_0\hookrightarrow\Lambda$\\
\bottomrule
\end{tabular}
\end{adjustbox}
\end{center}
Here $hH+\ell L$ is the exterior degree-two class, $c_0$ is an exterior
constant, and $(c_i)$ are filling constants.
The nonzero connecting maps are
\begin{equation}\label{eq:connecting-maps}
 \begin{aligned}
 \partial_Y^3(v)&=(\Sigma v)B,&
 \partial_Y^5(v)&=(\Sigma v)[\mathrm{pt}],\\
 \partial_{X_t}^2&:\Lambda\twoheadrightarrow W\hookrightarrow H^3(X_t),&
 \partial_{X_t}^5(v)&=(\Sigma v)[\mathrm{pt}].
 \end{aligned}
\end{equation}
All other connecting maps vanish. The signs follow by representing the
connecting class on an outward-$U$ collar by $d\chi_F\wedge\eta$,
where $\chi_F$ increases toward the filling. Pairing the degree-three
connecting class of $\beta_i$ with $H,L$ gives zero and one respectively,
hence $+B$. The same collar integral gives the top-degree sum maps.

Restriction from $Y$ preserves $H,L$, sends $(A,B)$ to $(\bar A,0)$,
and identifies $H^3(Y)$ with $\ker r_U$. Restriction from $X_t$ sends
$(H,a_4)$ to $(H,\bar A)$ and surjects $H^3(X_t)\to H^3(U)$ with
kernel the cohomological copy of $W$ in \eqref{eq:connecting-maps}.
Restriction to each Milnor filling is the corresponding component of
$r_U(y)$. These data specify the full integral cohomological attachment
diagram; no homotopy classification of the attaching maps is asserted.

\subsection{Monodromy and specialization}\label{sec:specialization}
Use left Poincar\'e duals in odd degree, characterized by
$\int\PD_L(\delta)\wedge\eta=\int_\delta\eta$, and orient $\delta_i$
so that $\PD_L(\delta_i)=\partial_{X_t}^2(\alpha_i)$.
For a positive simple loop, the Picard--Lefschetz transformation has
$N=T-I$, $N^2=0$, and, up to the common intersection-sign convention,
\begin{equation}\label{eq:monodromy}
 N\gamma=\pm\sum_{i=1}^{16}\langle\gamma,\delta_i\rangle\delta_i.
\end{equation}
Equivalently it factors through
$H_3(X_t)\to \Lambda_0\hookrightarrow\Lambda\twoheadrightarrow W$.
In the bases $e_i-e_{16}$ of $\Lambda_0$ and $\delta_1,\ldots,\delta_{15}$
of $W$, the nonzero block is $\pm(I_{15}+J_{15})$, where $J_{15}$
is the all-ones matrix. Its Smith form is
$\operatorname{diag}(1^{14},16)$: subtracting the first row from the
others gives fourteen unit pivots, and the remaining determinant is
sixteen. Thus the monodromy image has index sixteen in the primitive
vanishing lattice. This relation-matrix description is consistent with
the general conifold monodromy calculation of \cite{lee-lin-wang}.

The contraction $q:Y\to X_0$ collapses each exceptional sphere to its
node. Along the chosen smoothing ray, the Milnor collapse
$s_t:X_t\to X_0$ collapses each vanishing sphere. A disk cotangent bundle
of $S^3$ modulo its zero section is the cone on its sphere bundle; the
local collapses glue to the common exterior identification. Both maps
have degree one. Collapsing the sixteen cores to their distinct nodes
identifies the corresponding relative pairs with $(X_0,\{p_i\})$.
In degrees at least two its relative groups are $H_j(X_0)$. The pair
sequences give
\begin{equation}\label{eq:central-homology}
 \begin{split}
 &0\to H_3(Y)\xrightarrow{q_*}H_3(X_0)\to \Lambda_0\to0,\\
 &0\to W\to H_3(X_t)\xrightarrow{s_{t*}}H_3(X_0)\to0.
 \end{split}
\end{equation}
In degree two, $q_*$ is the quotient by $\ZZ c_Y$ and $s_{t*}$ is
an isomorphism. In degree four, $q_*$ is an isomorphism and $s_{t*}$
is a primitive injection with free rank-one cokernel. To see the latter,
the full relative sequence is
\[
 0\to H_4(X_t)\to H_4(X_0)\to\Lambda
 \xrightarrow{e_i\mapsto\delta_i}H_3(X_t)\to H_3(X_0)\to0,
\]
and the middle kernel is $\ZZ\mathbf1$. All central groups are free,
with Betti vector $(1,0,1,189,2,0,1)$. The maps are isomorphisms in
degrees zero and six; degrees one and five vanish.

The first sequence of \eqref{eq:exterior-homology} maps to the second
sequence of \eqref{eq:central-homology}, with identity on $W$, exterior
inclusion in the middle, and $q_*$ on the right. Both vertical cokernels
are canonically $\Lambda_0$. This is naturality of the actual collapse maps,
rather than an arbitrary splitting of free groups.

Let $h\in H^2(X_0)$ be the central polarization and choose
$A_0,B_0\in H^4(X_0)$ by $q^*(h,A_0,B_0)=(H,A,B)$.
The central even products and smoothing pullback are
\begin{equation}\label{eq:central-ring}
 h^2=5A_0+4B_0,\quad hA_0=[\mathrm{pt}],\quad hB_0=0,
 \qquad s_t^*(h,A_0,B_0)=(H,a_4,0).
\end{equation}
Indeed pairing the degree-four pullbacks with $H$, and using degree one,
fixes their coefficients integrally. Dually the image of $s_{t*}$ in
degree four is $\ZZ q_*\PD_Y(H)$. In degree three $q^*$ is surjective
with primitive kernel $\Lambda_0^*\simeq W$, and $s_t^*$ is a primitive injection
onto the annihilator of the vanishing homology lattice. Since $\im N$
has finite index in $W$, an integer-valued homomorphism annihilates
$\im N$ exactly when it annihilates $W$. Thus
\[
 \im s_t^*=\ker(T^*-I)\subset H^3(X_t;\ZZ).
\]
The central degree-three cup pairing factors through $q^*$ and the
unimodular pairing of $Y$, by naturality and degree one. Its radical is
$\ker q^*$; the Milnor pullback takes that radical to the primitive
vanishing cohomology lattice by naturality of the connecting map.
These formulas specify the central ring, including its degenerate
middle pairing, without assuming ordinary Poincar\'e duality for $X_0$.

